\subsection{Свойства квадратичной формы \texorpdfstring{\((Ax,x)\)}{(Ax,x)} и собственных значений самосопряженного оператора \texorpdfstring{\(A\)}{A}}

\begin{definition}[Самосопряженный оператор]
	Пусть \(H\) --- гильбертово пространство над \(\mathbb{C}\), и
	\[
	A\in\mathcal{L}(H).
	\]
	Оператор \(A\) называется \textbf{самосопряженным}, если
	\[
	A^*=A.
	\]
	Эквивалентно,
	\[
	(Ax,y)=(x,Ay)
	\qquad
	\forall x,y\in H.
	\]
\end{definition}

\begin{definition}[Квадратичная форма самосопряженного оператора]
	Пусть \(A\in\mathcal{L}(H)\). С оператором \(A\) связывают квадратичную форму
	\[
	K:H\to\mathbb{C},
	\]
	заданную правилом
	\[
	K(x)=(Ax,x).
	\]
	Идея состоит в том, что вместо оператора \(A\) часто удобно изучать значения формы
	\[
	(Ax,x).
	\]
\end{definition}

\begin{theorem}[Свойства квадратичной формы и собственных значений]
	Пусть \(H\) --- гильбертово пространство над \(\mathbb{C}\), и
	\[
	A\in\mathcal{L}(H)
	\]
	--- самосопряженный оператор. Тогда:
	
	\begin{itemize}
		\item для любого \(x\in H\)
		\[
		(Ax,x)\in\mathbb{R};
		\]
		
		\item если \(\lambda\) --- собственное значение оператора \(A\), то
		\[
		\lambda\in\mathbb{R};
		\]
		
		\item если \(\lambda_1\neq\lambda_2\) --- собственные значения оператора \(A\), а
		\(e_1,e_2\) --- соответствующие им собственные векторы, то
		\[
		(e_1,e_2)=0.
		\]
	\end{itemize}
\end{theorem}

\begin{proof}
	Докажем три утверждения по отдельности.
	
	\textbf{1. Значения квадратичной формы вещественны.}
	
	Пусть
	\[
	x\in H.
	\]
	Так как \(A\) самосопряжен, имеем
	\[
	(Ax,x)=(x,Ax).
	\]
	По свойству скалярного произведения
	\[
	(x,Ax)=\overline{(Ax,x)}.
	\]
	Следовательно,
	\[
	(Ax,x)=\overline{(Ax,x)}.
	\]
	Число совпадает со своим комплексным сопряжением, значит
	\[
	(Ax,x)\in\mathbb{R}.
	\]
	
	\textbf{2. Собственные значения самосопряженного оператора вещественны.}
	
	Пусть \(\lambda\) --- собственное значение оператора \(A\). Тогда существует ненулевой вектор
	\[
	e\in H,
	\qquad
	e\neq0,
	\]
	такой что
	\[
	Ae=\lambda e.
	\]
	Рассмотрим значение квадратичной формы на \(e\):
	\[
	(Ae,e).
	\]
	С одной стороны, по первому пункту
	\[
	(Ae,e)\in\mathbb{R}.
	\]
	С другой стороны,
	\[
	(Ae,e)
	=
	(\lambda e,e)
	=
	\lambda(e,e)
	=
	\lambda\|e\|^2.
	\]
	Так как \(e\neq0\), то
	\[
	\|e\|^2>0.
	\]
	Значит,
	\[
	\lambda
	=
	\frac{(Ae,e)}{\|e\|^2}
	\in\mathbb{R}.
	\]
	
	\textbf{3. Собственные векторы разных собственных значений ортогональны.}
	
	Пусть
	\[
	Ae_1=\lambda_1 e_1,
	\qquad
	Ae_2=\lambda_2 e_2,
	\]
	где
	\[
	\lambda_1\neq\lambda_2.
	\]
	Тогда
	\[
	(Ae_1,e_2)
	=
	(\lambda_1 e_1,e_2)
	=
	\lambda_1(e_1,e_2).
	\]
	С другой стороны, так как \(A\) самосопряжен,
	\[
	(Ae_1,e_2)
	=
	(e_1,Ae_2)
	=
	(e_1,\lambda_2 e_2).
	\]
	По пункту 2 собственные значения самосопряженного оператора вещественны, поэтому
	\[
	(e_1,\lambda_2 e_2)
	=
	\lambda_2(e_1,e_2).
	\]
	Значит,
	\[
	\lambda_1(e_1,e_2)
	=
	\lambda_2(e_1,e_2).
	\]
	Отсюда
	\[
	(\lambda_1-\lambda_2)(e_1,e_2)=0.
	\]
	Так как
	\[
	\lambda_1\neq\lambda_2,
	\]
	получаем
	\[
	(e_1,e_2)=0.
	\]
\end{proof}

\begin{remark}[Обратное утверждение]
	Верно и обратное: если для любого \(x\in H\)
	\[
	(Ax,x)\in\mathbb{R},
	\]
	то оператор \(A\) самосопряжен.
	
	Это доказывается поляризацией: рассматривают значения формы на \(x+y\) и \(x+iy\), после чего
	из вещественности соответствующих выражений получают
	\[
	(Ax,y)=(x,Ay).
	\]
	Для основного билета это утверждение обычно можно оставить как дополнительное.
\end{remark}